\section{Lean-Backed Semantic Target}

This section will present the Level B bridge implemented in
\code{SocialChoiceAtlas/Sen/RightAtomBridge.lean}. The bridge is semantic-only:
it backs the central M4 right-condition target, but it does not prove Python
clause correctness, CNF correctness, or artifact-level semantic-to-CNF
correctness.

The bridge defines:

\[
\texttt{RightAtomSemantics}\ F\ i\ x\ y :=
x \ne y \land \texttt{Decisive}\ F\ i\ x\ y.
\]

The inspected Lean file confirms orientation invariance through
\code{rightAtomSemantics_symm} and \code{rightAtomSemantics_iff_swap}, using
the relevant \code{Decisive.symm} result. It also confirms the sufficient
direction \code{twoRightAtoms_imply_MINLIB}: two fixed right conditions for
distinct voters imply Lean \code{MINLIB}.

TODO: Explain why this is enough to name a Lean-backed semantic target for the
central condition, but not enough to prove full selector-free equivalence to
\code{MINLIB}.

Safe claim-boundary note: say ``Lean-backed semantic target'' and ``sufficient
Lean bridge.'' Do not say that the Python clauses, CNF artifacts, or full
encoding have been proved correct.
